% =====================================================================
% Sazba zdrojového kódu
%
% Barvy jsou společné pro všechny jazyky, aby dokument s ukázkami ve
% více jazycích působil jednotně.
% =====================================================================

\definecolor{codeblue}{HTML}{007EED}
\definecolor{codegreen}{rgb}{0, 0.6, 0}
\definecolor{codegray}{rgb}{0.5, 0.5, 0.5}
\definecolor{codebeige}{HTML}{D4A000}
\definecolor{codeback}{rgb}{0.97, 0.97, 0.95}

\lstdefinestyle{latte}{
    backgroundcolor=\color{codeback},
    basicstyle=\footnotesize\ttfamily,
    keywordstyle=\color{codeblue}\bfseries,
    commentstyle=\color{codegreen},
    stringstyle=\color{codebeige},
    numberstyle=\tiny\color{codegray},
    identifierstyle=\color{black},
    numbers=left,
    numbersep=8pt,
    frame=single,
    rulecolor=\color{codegray!40},
    breaklines=true,
    breakatwhitespace=false,
    captionpos=b,
    showspaces=false,
    showstringspaces=false,
    showtabs=false,
    tabsize=4,
    extendedchars=true,
    xleftmargin=12pt,
    framexleftmargin=12pt,
}

\lstset{style=latte}

% Rust není v listings vestavěný, definice jazyka proto musí být vlastní.
\lstdefinelanguage{Rust}{
    keywords={as, async, await, break, const, continue, crate, dyn, else,
        enum, extern, false, fn, for, if, impl, in, let, loop, match, mod,
        move, mut, pub, ref, return, self, Self, static, struct, super,
        trait, true, type, unsafe, use, where, while},
    sensitive=true,
    morecomment=[l]{//},
    morecomment=[s]{/*}{*/},
    morestring=[b]",
}

% Krátký úryvek kódu uvnitř odstavce.
\newcommand{\inlinecode}[1]{\texttt{\colorbox{codeback}{#1}}}
